\subsection{Projective orbit identifications for a two-column band}

\begin{lemma}[The kernel of the constant-action quotient]
\label{C:two-column-projective-orbit-kernel}
Let $M$ have rectangular-band minimal ideal $K=I\times\{0,1\}$,
and let $S$ be a nonempty finite right $M$-set with no state fixed
by all of $K$.  Write $o$ for the number of right $K$-orbits in
$SK$, and let $\bar S$ be the constant-action quotient of
Lemma~\ref{C:constant-action-quotient-H2}.  Put
$R=\mathbb ZM$, $L=\widetilde{\mathbb ZS}$, and
$\bar L=\widetilde{\mathbb Z\bar S}$.  For any minimal idempotent
$e\in K$ there is an exact sequence of right $R$-modules
\[
 0\longrightarrow(eR)^{\oplus(o-1)}
   \longrightarrow L\longrightarrow\bar L\longrightarrow0.
\]
In particular, if $X=B\mathcal C(M,S)$ and
$Y=B\mathcal C(M,\bar S)$, their actual quotient map induces
isomorphisms $H_n(X;\mathbb Z)\cong H_n(Y;\mathbb Z)$ for
$n\geq3$, and an exact sequence
\[
 0\longrightarrow H_2(X;\mathbb Z)
 \longrightarrow H_2(Y;\mathbb Z)
 \longrightarrow\mathbb Z^{o-1}
 \longrightarrow H_1(X;\mathbb Z)
 \longrightarrow H_1(Y;\mathbb Z)\longrightarrow0.
\]
\end{lemma}
\begin{proof}
Every $K$-orbit in $SK$ has two elements: its column map has
at most two values, and a one-point orbit would be a $K$-fixed
state.  Label these orbits
\[
 O_a=\{s_{a,0},s_{a,1}\},\qquad 0\leq a<o,
\]
by their column values.  Each orbit is $M$-invariant.  Indeed,
if $x=y(i,j)$ and $m\in M$, then
$xm=y((i,j)m)=y(i,\tau_m(j))$ stays in the same orbit.
More precisely every actor has the same column map on every orbit:
\[
 s_{a,j}m=s_{a,\tau_m(j)}.
\]
The kernel of $L\to\bar L$ is $LK$, as proved in
Lemma~\ref{C:constant-action-quotient-H2}.  Its exact integral
basis is
\[
 s_{a,j}-s_{0,j}\qquad(1\leq a<o,\ j\in\{0,1\}).
\]
These differences belong to $LK$ by applying a band element in
column $j$ to points in the two indicated core orbits.
Conversely, for fixed $k=(i,j)$, every point of $Sk$ is the
column-$j$ point of one of these orbits.  Thus every generating
difference $sk-s'k$ is a difference of the displayed kind.
The points of $S\setminus SK$ introduce no additional relations.
The formula for the action identifies each pair indexed by $a$
with the right ideal $eR$, whose basis is the two elements in
the row of $e$.  This proves the module statement, including
projectivity of its kernel.

The proof of Lemma~\ref{C:act-cone-Tor} works over $\mathbb Z$
without change: use a projective left $\mathbb ZM$-resolution
of the trivial module and the exact augmentation sequence of
the free abelian permutation module $\mathbb ZS$.  Thus
$\widetilde H_n(X;\mathbb Z)=\operatorname{Tor}_{n-1}^R(L,\mathbb Z)$
for every $n\geq1$, naturally in $S$.  Apply the Tor exact
sequence to the displayed projective-kernel sequence and note
that $eR\otimes_R\mathbb Z\cong\mathbb Z$.  This gives exactly
the asserted homology sequence and all higher isomorphisms.
\end{proof}

\begin{lemma}[A free second-homology complement]
\label{C:free-second-homology-wedge}
Let $f:X\to Y$ be a map of simply connected CW complexes.
Suppose $f_*$ is an isomorphism on integral homology in every
degree at least three, and fits into an exact sequence
\[
 0\longrightarrow H_2(X;\mathbb Z)
 \xrightarrow{f_*}H_2(Y;\mathbb Z)
 \longrightarrow\mathbb Z^r\longrightarrow0.
\]
There is a homotopy equivalence
$X\vee\bigvee_rS^2\to Y$ whose restriction to $X$ is $f$,
up to homotopy. No finiteness assumption on $X$ or $Y$ is needed.
\end{lemma}
\begin{proof}
The displayed sequence splits because its quotient is free.
Choose a basis of a complementary copy of $\mathbb Z^r$ in
$H_2(Y;\mathbb Z)$. Since $Y$ is simply connected, the Hurewicz
isomorphism $\pi_2(Y)\cong H_2(Y;\mathbb Z)$ represents this
basis by based maps from $r$ two-spheres. Wedge these maps with
$f$, after the harmless choice of compatible basepoints.
The resulting map is an isomorphism on second homology by the
chosen splitting, and on every higher homology group by the
hypothesis on $f$. It is also an isomorphism in degrees zero
and one. Both spaces are simply connected CW complexes, so
homological Whitehead makes it a homotopy equivalence.
\end{proof}

\begin{corollary}[The orbit quotient splits off two-spheres]
\label{C:two-column-orbit-wedge}
Under the hypotheses of
Lemma~\ref{C:two-column-projective-orbit-kernel}, suppose the
original act cone $X=B\mathcal C(M,S)$ is simply connected.
The actual quotient map to
$Y=B\mathcal C(M,\bar S)$ extends to a homotopy equivalence
\[
 X\vee\bigvee_{o-1}S^2\ \simeq\ Y\ \simeq\ BN,
\]
where $N$ is the constant-action compression monoid.
\end{corollary}
\begin{proof}
The compression theorem gives the displayed equivalence with
$BN$. Its minimal ideal is a rectangular band, so its universal
group and hence its fundamental group are trivial.
The projective-kernel lemma gives the required isomorphisms in
all degrees at least three. Since $H_1(X)=H_1(Y)=0$, its exact
sequence in low degrees has free cokernel $\mathbb Z^{o-1}$.
Apply Lemma~\ref{C:free-second-homology-wedge} to the actual
quotient map. The equivalence thus retains that map on the
original summand; it is not merely an abstract homology comparison.
\end{proof}

\begin{lemma}[A finite model for the quotient of a hypothetical projective plane]
\label{C:two-column-finite-quotient-model}
Under the hypotheses of
Lemma~\ref{C:two-column-projective-orbit-kernel}, suppose
$B\mathcal C(M,S)\simeq\mathbb CP^2$.
The compressed monoid $N$ associated to $\bar S$ has the same
number of outside elements as $M$, minimal band
$I\times\bar S$, and a finite four-dimensional CW model with
one zero-cell, $o$ two-cells, and one four-cell. More precisely,
\[
 BN\simeq\mathbb CP^2\vee\bigvee_{o-1}S^2.
\]
Its integral
homology is
\[
 H_2(BN)=\mathbb Z^o,\qquad H_4(BN)=\mathbb Z,
 \qquad \widetilde H_n(BN)=0\quad(n\notin\{2,4\}).
\]
The map from the original act cone to $BN$ is surjective on
integral degree-two cohomology in the pullback direction.
\end{lemma}
\begin{proof}
Compression is Theorem~\ref{C:constant-minimal-act-compression};
it leaves the outside elements unchanged and replaces the band
by $I\times\bar S$.  The preceding exact sequence gives the
homology assertion, since an extension of $\mathbb Z^{o-1}$
by $\mathbb Z$ splits as an abelian group.
Corollary~\ref{C:two-column-orbit-wedge} proves the stronger
wedge identification under the original projective-plane summand.
The usual cell structure of $\mathbb CP^2$, together with the
additional two-spheres, is the stated finite four-dimensional model.
Finally the asserted cohomology surjection is the integral
statement of Lemma~\ref{C:constant-action-quotient-H2} followed
by the compression equivalence.
\end{proof}

\begin{theorem}[Four outside elements never suffice for a two-column act cone]
\label{C:two-column-four-outside-act-cones}
Let $M$ have a rectangular-band minimal ideal with two columns
and at most four elements outside it.  For every nonempty finite
right $M$-set $S$,
\[
 B\mathcal C(M,S)\not\simeq\mathbb CP^2.
\]
There is no bound on the number of rows or states.
More generally, if the act cone has bounded rational homology, all
of its rational degree-two products vanish.  Thus the obstruction
also applies to the rational homology and cohomology ring of a
projective plane, without requiring an integral finite CW model.
\end{theorem}
\begin{proof}
We prove the stronger rational assertion.
If $S$ has a state fixed by $K$, the band-grid argument in
Lemma~\ref{C:four-state-band-dichotomy} makes every degree-two
class relative to $BM$, and their products are zero.  That argument
does not use the number of states.
Otherwise Lemma~\ref{C:two-column-projective-orbit-kernel} and
compression give a monoid $N$ with the same rational homology
as the act cone in every degree at least three.  Its rational
homology is therefore bounded.  It has at most four elements
outside its minimal band.  Theorem~\ref{B:contracted-four-Tor},
applied to its row and column augmentation modules in the band
Tor formula, implies $H_4(BN;\mathbb Q)=0$: a nonzero group in
that degree would force nonzero groups in arbitrarily high degrees.
Thus the original act cone also has zero rational fourth homology,
and its rational degree-two products vanish.  This argument uses
neither an integral finiteness inference nor a degree-two rank
restriction on the compressed monoid.
\end{proof}

\begin{corollary}[Six and seven nonconstant states over three rows]
\label{C:three-row-six-seven-nonconstant}
Let $M$ have minimal band $3\times2$, and let $S$ have six or
seven elements.  If some element of the band acts nonconstantly
on $S$, then $B\mathcal C(M,S)\not\simeq\mathbb CP^2$.
There is no bound on the order of $M$.
\end{corollary}
\begin{proof}
Suppose the cone is a projective plane.  As above, it has no
$K$-fixed state.  The act is connected: the homotopy-pushout
homology sequence identifies its first rational homology with
the reduced degree-zero homology of the action category, because
$H_1(BM;\mathbb Q)=0$.
If $SK$ had a single orbit, every band actor would be constant
on all states, since $xk=(xk)k$ and $k$ is constant on that orbit.
Thus $o\geq2$.  The $K$-orbits are individually $M$-invariant,
so $S=SK$ with $o>1$ would disconnect the act.

For six states these facts force $o=2$, with two states outside
$SK$.  The constant-action quotient identifies the two core
orbits column by column and leaves the other states distinct,
so $|\bar S|=4$.  The finite compressed model has minimal band
$3\times4$ and degree-two rank two.  The rank-two portion of
the proof of Theorem~\ref{B:three-by-four-targets} makes every
rational degree-two product zero for any finite CW source of
that rank: its only nonzero-square exceptional orbit has rank
one.  Its actual zero-product certificates and local-group
exclusions therefore apply to this compressed model.

For seven states, either $o=2$ and $|\bar S|=5$, or $o=3$ and
$|\bar S|=3$.  In the former case the finite compressed model
has degree-two rank two and minimal band $3\times5$; every
rank-two orbit in the proof of
Theorem~\ref{B:three-by-five-targets} has zero rational products.
In the latter case Theorem~\ref{B:three-by-three-cups} applies
without a degree-two rank restriction.
In all cases every rational degree-two class of the original
cone pulls back from the compressed model.  Its nonzero square
is therefore impossible.
\end{proof}

\begin{remark}[Scope of the orbit argument]
\label{C:two-column-orbit-quotient-scope}
The projective-kernel assertion uses a minimal band with exactly
two columns and the absence of one-point core orbits.  For a
larger column set, different core orbits can identify different
columns, and the displayed kernel need not be projective.
The proof does not assert projectivity in that situation.
The integral orbit calculation and the finite four-cell model
were independently audited by agents B and D; the six-state
rank-two application was checked against the full earlier
classification, not inferred merely from its named product target.
\end{remark}
